\chapter{Smart Contract Configuration Parameters}\label{ch:appendix-a-parameters}
This appendix provides a comprehensive reference for all configurable parameters accepted by the smart contract at deployment time.
The constructor receives a
    \texttt{Config} struct, which is composed of four sub-structures:
        \texttt{SmartContractNeeds},
        \texttt{ChunkConfig},
        \texttt{RewardConfig} and \texttt{GlobalPosition}.
Each sub-structure is documented in a dedicated table below.
In addition, the constructor receives the \texttt{B} public key and the address of \texttt{DataSentTokenContract} which is the token introduced for senders and verifiers rewarding.
The details of rewarding model are described in the~\ref{subsec:proto-rewarding-model} subsection.

A convention adopted throughout the configuration is the \textbf{zero-means-unlimited} semantics: when a parameter that governs a limit is set to zero, the corresponding constraint is not enforced.
This convention is explicitly noted for each applicable parameter.

% SmartContractNeeds
\section{Campaign Requirements (\texttt{SmartContractNeeds})}\label{sec:appendix-a-needs}
The \texttt{SmartContractNeeds} structure defines the operational rules of the crowdsensing campaign, including participation limits, and data acceptance criteria.
\ref{tab:params-needs} lists each parameter with its type, constraints, and description.
\begin{table}
    \centering
    \caption{Parameters of the \texttt{SmartContractNeeds} structure.}
    \label{tab:params-needs}
    \small
    \begin{tabular}{p{3.8cm} l p{1.8cm} p{5.8cm}}
        \hline
        \textbf{Parameter} & \textbf{Type} & \textbf{Constraint} & \textbf{Description} \\
        \hline
        %\texttt{freshnessPeriod}
        %& \texttt{uint256}
        %& $\geq 0$ (0 = disabled)
        %& Reserved for future use. Intended to define a time window for data freshness validation. Currently not enforced by the contract logic. \\[6pt]
        \texttt{dataNeeded}
        & \texttt{uint256}
        & $\geq 0$ (0 = unlimited)
        & Number of positively verified data entries required to automatically close the campaign.
        When set to zero, the campaign does not auto-close based on data count alone. \\[6pt]

        \texttt{maximumDataPerChunk}
        & \texttt{uint256}
        & $\geq 0$ (0 = unlimited)
        & Maximum number of data submissions accepted within a single spatial chunk.
        Prevents over-representation of a geographic area. \\[6pt]

        \texttt{checksRequired}
        & \texttt{uint8}
        & odd, $\geq 1$
        & Number of independent verification slots allocated per data entry.
        The odd constraint ensures that a strict majority can always be determined. \\[6pt]

        \texttt{verifierCooldown}
        & \texttt{uint256}
        & $> 0$
        & Minimum distance (in data entry indices) between consecutive assignments of the same verifier.
        Prevents a single verifier from being assigned to sequentially adjacent data entries, promoting verification independence. \\[6pt]

        \texttt{verificationTimeout}
        & \texttt{uint256}
        & $> 0$
        & Time window (in seconds) granted to a verifier to complete the assigned verification task.
        If the deadline expires, the verification slot becomes available for reassignment. \\[6pt]

        \texttt{minimumParticipants}
        & \texttt{uint256}
        & $\geq 0$ (0 = unset)
        & Minimum number of unique senders that must have at least one positively verified data entry for the campaign to auto-close.
        Combined with \texttt{dataNeeded} in a dual-threshold check. \\[6pt]

        \texttt{maxDataPerSender}
        & \texttt{uint256}
        & $\geq 0$ (0 = unlimited)
        & Maximum number of data entries a single sender is allowed to submit.
        When greater than zero, the bond mechanism is activated: senders must deposit a bond proportional to this limit before submitting data. \\[6pt]

        \texttt{maxDataPerVerifier\-PerBond}
        & \texttt{uint256}
        & $\geq 0$
        & Number of verification tasks covered by a single bond deposit.
        Used to compute the required verifier bond as $\texttt{maxDataPerVerifierPerBond} \times \texttt{dataVerifierBondPerData}$. \\
        \hline
    \end{tabular}
\end{table}


% ChunkConfig
\section{Spatial Chunking (\texttt{ChunkConfig})}\label{sec:appendix-a-chunks}

The \texttt{ChunkConfig} structure controls the spatial partitioning of the campaign area into a square grid of chunks.
\cref{tab:params-chunks} lists the two parameters involved.

\begin{table}
    \centering
    \caption{Parameters of the \texttt{ChunkConfig} structure.}
    \label{tab:params-chunks}
    \small
    \begin{tabular}{p{3.8cm} l p{2.8cm} p{4.8cm}}
        \hline
        \textbf{Parameter} & \textbf{Type} & \textbf{Constraint} & \textbf{Description} \\
        \hline
        \texttt{chunkSize}
        & \texttt{uint256}
        & $> 0$
        & Side length of each square chunk, expressed in E7 coordinate units (i.e., $10^{-7}$ degrees).
        For reference, 1\,000 E7 units correspond to approximately 11 meters at the equator. \\[6pt]

        \texttt{chunk4Side}
        & \texttt{uint16}
        & $> 0$
        & Number of chunks along each side of the square campaign area.
        The total number of chunks is $\texttt{chunk4Side}^2$.

        When set to 1, the spatial chunking system is effectively disabled and the entire campaign area is treated as a single region. \\
        \hline
    \end{tabular}
\end{table}

The total geographic extent of the campaign area is derived as $\texttt{chunkSize} \times \texttt{chunk4Side}$ E7 units per side, starting from the origin coordinates specified in the \texttt{GlobalPosition} parameter (see \Cref{sec:appendix-a-origin-coordinated}.
% RewardConfig
\section{Reward and Bond Configuration (\texttt{RewardConfig})}\label{sec:appendix-a-rewards}

The \texttt{RewardConfig} structure governs the economic model of the campaign, defining the total token supply, per-operation payouts, and bond amounts denominated in DST (DataSent Token).
\ref{tab:params-rewards} lists all parameters.
\begin{table}
    \centering
    \caption{Parameters of the \texttt{RewardConfig} structure.}
    \label{tab:params-rewards}
    \small
    \begin{tabular}{p{3.8cm} l p{2.8cm} p{4.8cm}}
        \hline
        \textbf{Parameter} & \textbf{Type} & \textbf{Constraint} & \textbf{Description} \\
        \hline
        \texttt{totalDst}
        & \texttt{uint256}
        & $> 0$
        & Total amount of DST tokens that the campaign initiator must deposit into the contract to fund all sender and verifier rewards. \\[6pt]

        \texttt{dataSenderPayout}
        & \texttt{uint16}
        & $>$ \texttt{dataSender\-BondPerData}
        & DST tokens paid to a sender when one of their data entries is positively verified.
        The constraint ensures that participation is always economically incentivized. \\[6pt]

        \texttt{dataVerifierPayout}
        & \texttt{uint16}
        & $>$ \texttt{dataVerifier\-BondPerData}
        & DST tokens paid to a verifier who voted with the majority on a resolved data entry. \\[6pt]

        \texttt{dataSenderBond\-PerData}
        & \texttt{uint16}
        & $> 0$
        & DST tokens deducted from the sender's bond balance at data submission time.
        Returned together with the payout upon positive verification. \\[6pt]

        \texttt{dataVerifierBond\-PerData}
        & \texttt{uint16}
        & $> 0$
        & DST tokens deducted from the verifier's bond balance when a verification task is assigned.
        Returned together with the payout upon a majority-aligned vote. \\
        \hline
    \end{tabular}
\end{table}

The bond mechanism serves a dual purpose: it discourages low-quality submissions and frivolous verifications, while ensuring that participants have a stake in the correctness of their contributions.
Bonds are deducted at action time (submission or assignment) rather than at reward time, so that the balance check is enforced before the operation is recorded on-chain.


% GlobalPosition
\section{Campaign Origin Coordinates (\texttt{GlobalPosition})}\label{sec:appendix-a-origin-coordinated}

The \texttt{GlobalPosition} structure defines the geographic origin (south-west corner) of the campaign area.
The coordinate representation and the E7 packing mechanism are described in detail in the implementation chapter. \Cref{tab:params-position} lists the fields and their valid ranges.

\begin{table}[H]
    \centering
    \caption{Fields of the \texttt{GlobalPosition} structure.}
    \label{tab:params-position}
    \small
    \begin{tabular}{p{3.8cm} l p{2.8cm} p{4.8cm}}
        \hline
        \textbf{Field} & \textbf{Type} & \textbf{Valid Range} & \textbf{Description} \\
        \hline
        \texttt{latitudeInteger}
        & \texttt{int32}
        & $[-90,\; +90]$
        & Integer part of the latitude in degrees. \\[6pt]

        \texttt{latitudeDecimal}
        & \texttt{uint32}
        & $[0,\; 9\,999\,999]$
        & Fractional part of the latitude, representing up to seven decimal digits. \\[6pt]

        \texttt{longitudeInteger}
        & \texttt{int32}
        & $[-180,\; +180]$
        & Integer part of the longitude in degrees. \\[6pt]

        \texttt{longitudeDecimal}
        & \texttt{uint32}
        & $[0,\; 9\,999\,999]$
        & Fractional part of the longitude, representing up to seven decimal digits. \\
        \hline
    \end{tabular}
\end{table}